\block{Introducción}{
	Puerto Rico es la jurisdicción de los Estados Unidos que consistentemente presenta la mayor brecha entre el Producto Interno Bruto (PIB)
	y el Producto Nacional Bruto (PNB) en comparación con otras economías avanzadas \cite{torres2021puerto}.
	Esta divergencia estructural, impulsada en gran medida por la repatriación de ganancias de corporaciones
	no residentes, dificulta la medición precisa de la actividad productiva real a nivel regional.

	\vspace{0.5em}
	Esta problemática ha motivado a entidades como el (BEA) a presentar
	e implementar reformas metodológicas para modernizar la producción de estadísticas macroeconómicas \cite{bea2011puertorico}.
	Varios estudios abarcan sobre la producción de PR pero pocos se enfocan en la producción de los municipios.
	Ante esta limitación, evaluar la eficiencia de los 78 municipios requiere de enfoques
	alternativos. Este trabajo aborda dicha brecha
}

